% Options for packages loaded elsewhere
\PassOptionsToPackage{unicode}{hyperref}
\PassOptionsToPackage{hyphens}{url}
%
\documentclass[
  ignorenonframetext,
  aspectratio=169,
  russian,
]{beamer}
\usepackage{pgfpages}
\setbeamertemplate{caption}[numbered]
\setbeamertemplate{caption label separator}{: }
\setbeamercolor{caption name}{fg=normal text.fg}
\beamertemplatenavigationsymbolshorizontal
% Prevent slide breaks in the middle of a paragraph
\widowpenalties 1 10000
\raggedbottom
\setbeamertemplate{part page}{
  \centering
  \begin{beamercolorbox}[sep=16pt,center]{part title}
    \usebeamerfont{part title}\insertpart\par
  \end{beamercolorbox}
}
\setbeamertemplate{section page}{
  \centering
  \begin{beamercolorbox}[sep=12pt,center]{section title}
    \usebeamerfont{section title}\insertsection\par
  \end{beamercolorbox}
}
\setbeamertemplate{subsection page}{
  \centering
  \begin{beamercolorbox}[sep=8pt,center]{subsection title}
    \usebeamerfont{subsection title}\insertsubsection\par
  \end{beamercolorbox}
}
\AtBeginPart{
  \frame{\partpage}
}
\AtBeginSection{
  \ifbibliography
  \else
    \frame{\sectionpage}
  \fi
}
\AtBeginSubsection{
  \frame{\subsectionpage}
}

\usepackage{amsmath,amssymb}
\usepackage{iftex}
\ifPDFTeX
  \usepackage[T1]{fontenc}
  \usepackage[utf8]{inputenc}
  \usepackage{textcomp} % provide euro and other symbols
\else % if luatex or xetex
  \usepackage{unicode-math}
  \defaultfontfeatures{Scale=MatchLowercase}
  \defaultfontfeatures[\rmfamily]{Ligatures=TeX,Scale=1}
\fi
\usepackage{lmodern}
\usetheme[]{Madrid}
\ifPDFTeX\else  
    % xetex/luatex font selection
\fi
% Use upquote if available, for straight quotes in verbatim environments
\IfFileExists{upquote.sty}{\usepackage{upquote}}{}
\IfFileExists{microtype.sty}{% use microtype if available
  \usepackage[]{microtype}
  \UseMicrotypeSet[protrusion]{basicmath} % disable protrusion for tt fonts
}{}
\usepackage{xcolor}
\newif\ifbibliography
\setlength{\emergencystretch}{3em} % prevent overfull lines
\setcounter{secnumdepth}{-\maxdimen} % remove section numbering


\providecommand{\tightlist}{%
  \setlength{\itemsep}{0pt}\setlength{\parskip}{0pt}}\usepackage{longtable,booktabs,array}
\usepackage{calc} % for calculating minipage widths
\usepackage{caption}
% Make caption package work with longtable
\makeatletter
\def\fnum@table{\tablename~\thetable}
\makeatother
\usepackage{graphicx}
\makeatletter
\newsavebox\pandoc@box
\newcommand*\pandocbounded[1]{% scales image to fit in text height/width
  \sbox\pandoc@box{#1}%
  \Gscale@div\@tempa{\textheight}{\dimexpr\ht\pandoc@box+\dp\pandoc@box\relax}%
  \Gscale@div\@tempb{\linewidth}{\wd\pandoc@box}%
  \ifdim\@tempb\p@<\@tempa\p@\let\@tempa\@tempb\fi% select the smaller of both
  \ifdim\@tempa\p@<\p@\scalebox{\@tempa}{\usebox\pandoc@box}%
  \else\usebox{\pandoc@box}%
  \fi%
}
% Set default figure placement to htbp
\def\fps@figure{htbp}
\makeatother

\usepackage{polyglossia}
\setdefaultlanguage{russian}
\setotherlanguage{english}
\usepackage{fontspec}
\setmainfont{FreeSerif}
\usepackage{polyglossia}
\setdefaultlanguage{russian}
\setotherlanguage{english}
\makeatletter
\@ifpackageloaded{caption}{}{\usepackage{caption}}
\AtBeginDocument{%
\ifdefined\contentsname
  \renewcommand*\contentsname{Содержание}
\else
  \newcommand\contentsname{Содержание}
\fi
\ifdefined\listfigurename
  \renewcommand*\listfigurename{Список иллюстраций}
\else
  \newcommand\listfigurename{Список иллюстраций}
\fi
\ifdefined\listtablename
  \renewcommand*\listtablename{Список таблиц}
\else
  \newcommand\listtablename{Список таблиц}
\fi
\ifdefined\figurename
  \renewcommand*\figurename{Рисунок}
\else
  \newcommand\figurename{Рисунок}
\fi
\ifdefined\tablename
  \renewcommand*\tablename{Таблица}
\else
  \newcommand\tablename{Таблица}
\fi
}
\@ifpackageloaded{float}{}{\usepackage{float}}
\floatstyle{ruled}
\@ifundefined{c@chapter}{\newfloat{codelisting}{h}{lop}}{\newfloat{codelisting}{h}{lop}[chapter]}
\floatname{codelisting}{Список}
\newcommand*\listoflistings{\listof{codelisting}{Листинги}}
\makeatother
\makeatletter
\makeatother
\makeatletter
\@ifpackageloaded{caption}{}{\usepackage{caption}}
\@ifpackageloaded{subcaption}{}{\usepackage{subcaption}}
\makeatother

\ifLuaTeX
\usepackage[bidi=basic]{babel}
\else
\usepackage[bidi=default]{babel}
\fi
\babelprovide[main,import]{russian}
% get rid of language-specific shorthands (see #6817):
\let\LanguageShortHands\languageshorthands
\def\languageshorthands#1{}
\usepackage{csquotes}
\usepackage{bookmark}

\IfFileExists{xurl.sty}{\usepackage{xurl}}{} % add URL line breaks if available
\urlstyle{same} % disable monospaced font for URLs
\hypersetup{
  pdftitle={Лабораторная работа №2},
  pdfauthor={София Черная},
  pdflang={ru-RU},
  hidelinks,
  pdfcreator={LaTeX via pandoc}}


\title{Лабораторная работа №2}
\subtitle{Модель SIR и модель Лотки-Вольтерры}
\author{София Черная}
\date{2026-03-07}

\begin{document}
\frame{\titlepage}

\renewcommand*\contentsname{Содержание}
\begin{frame}[allowframebreaks]
  \frametitle{Содержание}
  \setcounter{tocdepth}{1}
  \tableofcontents
\end{frame}

\section{Информация}\label{ux438ux43dux444ux43eux440ux43cux430ux446ux438ux44f}

\begin{frame}{Докладчик}
\phantomsection\label{ux434ux43eux43aux43bux430ux434ux447ux438ux43a}
\begin{columns}[c]
\begin{column}{0.7\linewidth}
\begin{itemize}[<+->]
\tightlist
\item
  София Черная
\item
  Студентка
\item
  Российский университет дружбы народов
\item
  \href{mailto:1132236043@pfur.ru}{\nolinkurl{1132236043@pfur.ru}}
\end{itemize}
\end{column}

\begin{column}{0.3\linewidth}
\end{column}
\end{columns}
\end{frame}

\section{Вводная
часть}\label{ux432ux432ux43eux434ux43dux430ux44f-ux447ux430ux441ux442ux44c}

\begin{frame}{Актуальность}
\phantomsection\label{ux430ux43aux442ux443ux430ux43bux44cux43dux43eux441ux442ux44c}
\begin{itemize}[<+->]
\tightlist
\item
  Моделирование эпидемий позволяет прогнозировать развитие заболеваний
\item
  Модели «хищник-жертва» важны для понимания экологических систем
\item
  Имитационное моделирование --- современный инструмент исследования
\item
  Julia и DrWatson --- современный стек для воспроизводимых исследований
\end{itemize}
\end{frame}

\begin{frame}{Цели и задачи}
\phantomsection\label{ux446ux435ux43bux438-ux438-ux437ux430ux434ux430ux447ux438}
\begin{itemize}[<+->]
\tightlist
\item
  Реализовать модель SIR в базовом варианте
\item
  Реализовать модель Лотки-Вольтерры
\item
  Преобразовать код в литературный стиль с помощью Literate.jl
\item
  Сгенерировать Jupyter Notebook и документацию Quarto
\item
  Провести параметрические исследования моделей
\item
  Проанализировать результаты и визуализировать графики
\end{itemize}
\end{frame}

\begin{frame}{Материалы и методы}
\phantomsection\label{ux43cux430ux442ux435ux440ux438ux430ux43bux44b-ux438-ux43cux435ux442ux43eux434ux44b}
\begin{itemize}[<+->]
\tightlist
\item
  Язык программирования Julia
\item
  Пакет DifferentialEquations.jl для решения ОДУ
\item
  Пакет DrWatson.jl для управления проектами
\item
  Пакет Literate.jl для литературного программирования
\item
  Пакет Plots.jl для визуализации
\item
  Пакет DataFrames.jl для анализа данных
\end{itemize}
\end{frame}

\section{Теоретическое
введение}\label{ux442ux435ux43eux440ux435ux442ux438ux447ux435ux441ux43aux43eux435-ux432ux432ux435ux434ux435ux43dux438ux435}

\begin{frame}{Модель SIR}
\phantomsection\label{ux43cux43eux434ux435ux43bux44c-sir}
\begin{itemize}[<+->]
\tightlist
\item
  Классическая эпидемиологическая модель
\item
  Популяция делится на три группы:

  \begin{itemize}[<+->]
  \tightlist
  \item
    \textbf{S (Susceptible)} --- восприимчивые
  \item
    \textbf{I (Infectious)} --- инфицированные
  \item
    \textbf{R (Recovered)} --- выздоровевшие
  \end{itemize}
\end{itemize}
\end{frame}

\begin{frame}{Система уравнений SIR}
\phantomsection\label{ux441ux438ux441ux442ux435ux43cux430-ux443ux440ux430ux432ux43dux435ux43dux438ux439-sir}
\[
\begin{cases}
\dfrac{dS}{dt} = -\beta \dfrac{S I}{N} \\[1em]
\dfrac{dI}{dt} = \beta \dfrac{S I}{N} - \gamma I \\[1em]
\dfrac{dR}{dt} = \gamma I
\end{cases}
\]

где: - ( N = S + I + R ) --- общая численность популяции - ( \beta ) ---
коэффициент заражения - ( \gamma ) --- коэффициент выздоровления
\end{frame}

\begin{frame}{Ключевые показатели}
\phantomsection\label{ux43aux43bux44eux447ux435ux432ux44bux435-ux43fux43eux43aux430ux437ux430ux442ux435ux43bux438}
\begin{itemize}[<+->]
\tightlist
\item
  \textbf{Базовое репродуктивное число}: ( R\_0 = \dfrac{\beta}{\gamma}
  )

  \begin{itemize}[<+->]
  \tightlist
  \item
    ( R\_0 \textgreater{} 1 ) --- эпидемия разрастается
  \item
    ( R\_0 \textless{} 1 ) --- эпидемия затухает
  \end{itemize}
\item
  \textbf{Порог коллективного иммунитета}: ( 1 - \dfrac{1}{R_0} )
\item
  \textbf{Эффективное репродуктивное число}: ( R\_e(t) = R\_0
  \cdot \frac{S(t)}{N} )
\end{itemize}
\end{frame}

\begin{frame}{Модель Лотки-Вольтерры}
\phantomsection\label{ux43cux43eux434ux435ux43bux44c-ux43bux43eux442ux43aux438-ux432ux43eux43bux44cux442ux435ux440ux440ux44b}
\begin{itemize}[<+->]
\tightlist
\item
  Модель «хищник-жертва»
\item
  Описывает взаимодействие двух биологических видов
\end{itemize}

\begin{block}{Система уравнений}
\phantomsection\label{ux441ux438ux441ux442ux435ux43cux430-ux443ux440ux430ux432ux43dux435ux43dux438ux439}
\[
\begin{cases}
\dfrac{dx}{dt} = \alpha x - \beta x y \\[1em]
\dfrac{dy}{dt} = \delta x y - \gamma y
\end{cases}
\]
\end{block}
\end{frame}

\begin{frame}{Параметры модели Лотки-Вольтерры}
\phantomsection\label{ux43fux430ux440ux430ux43cux435ux442ux440ux44b-ux43cux43eux434ux435ux43bux438-ux43bux43eux442ux43aux438-ux432ux43eux43bux44cux442ux435ux440ux440ux44b}
\begin{itemize}[<+->]
\tightlist
\item
  ( x ) --- численность жертв
\item
  ( y ) --- численность хищников
\item
  ( \alpha ) --- прирост жертв
\item
  ( \beta ) --- смертность жертв от хищников
\item
  ( \delta ) --- прирост хищников от жертв
\item
  ( \gamma ) --- смертность хищников
\end{itemize}
\end{frame}

\begin{frame}{Стационарная точка}
\phantomsection\label{ux441ux442ux430ux446ux438ux43eux43dux430ux440ux43dux430ux44f-ux442ux43eux447ux43aux430}
\begin{itemize}[<+->]
\tightlist
\item
  Нетривиальная точка сосуществования видов:
\end{itemize}

{[} x\^{}* = \frac{\gamma}{\delta},\quad y\^{}* = \frac{\alpha}{\beta}
{]}

\begin{itemize}[<+->]
\tightlist
\item
  Вокруг этой точки решения носят колебательный характер
\item
  Пик численности хищников отстаёт от пика численности жертв
\end{itemize}
\end{frame}

\section{Выполнение лабораторной
работы}\label{ux432ux44bux43fux43eux43bux43dux435ux43dux438ux435-ux43bux430ux431ux43eux440ux430ux442ux43eux440ux43dux43eux439-ux440ux430ux431ux43eux442ux44b}

\begin{frame}{Модель SIR}
\phantomsection\label{ux43cux43eux434ux435ux43bux44c-sir-1}
\begin{block}{Создание файла}
\phantomsection\label{ux441ux43eux437ux434ux430ux43dux438ux435-ux444ux430ux439ux43bux430}
\begin{figure}[H]

{\centering \includegraphics[width=0.7\linewidth,height=\textheight,keepaspectratio]{image/1.png}

}

\caption{Создание файла через команду touch}

\end{figure}%
\end{block}

\begin{block}{Программный код модели SIR}
\phantomsection\label{ux43fux440ux43eux433ux440ux430ux43cux43cux43dux44bux439-ux43aux43eux434-ux43cux43eux434ux435ux43bux438-sir}
\begin{figure}[H]

{\centering \includegraphics[width=0.7\linewidth,height=\textheight,keepaspectratio]{image/2.png}

}

\caption{Программный код модели SIR}

\end{figure}%
\end{block}

\begin{block}{Запуск скрипта}
\phantomsection\label{ux437ux430ux43fux443ux441ux43a-ux441ux43aux440ux438ux43fux442ux430}
\begin{figure}[H]

{\centering \includegraphics[width=0.7\linewidth,height=\textheight,keepaspectratio]{image/3.png}

}

\caption{Запуск скрипта sir\_ode.jl}

\end{figure}%
\end{block}

\begin{block}{Параметры модели}
\phantomsection\label{ux43fux430ux440ux430ux43cux435ux442ux440ux44b-ux43cux43eux434ux435ux43bux438}
\begin{figure}[H]

{\centering \includegraphics[width=0.7\linewidth,height=\textheight,keepaspectratio]{image/4.png}

}

\caption{Параметры модели SIR, выведенные в консоль}

\end{figure}%
\end{block}
\end{frame}

\begin{frame}{Результаты модели SIR}
\phantomsection\label{ux440ux435ux437ux443ux43bux44cux442ux430ux442ux44b-ux43cux43eux434ux435ux43bux438-sir}
\begin{block}{Динамика эффективного репродуктивного числа}
\phantomsection\label{ux434ux438ux43dux430ux43cux438ux43aux430-ux44dux444ux444ux435ux43aux442ux438ux432ux43dux43eux433ux43e-ux440ux435ux43fux440ux43eux434ux443ux43aux442ux438ux432ux43dux43eux433ux43e-ux447ux438ux441ux43bux430}
\begin{figure}[H]

{\centering \includegraphics[width=0.7\linewidth,height=\textheight,keepaspectratio]{image/5.png}

}

\caption{Динамика эффективного репродуктивного числа ( R\_e(t) )}

\end{figure}%
\end{block}

\begin{block}{Динамика числа зараженных}
\phantomsection\label{ux434ux438ux43dux430ux43cux438ux43aux430-ux447ux438ux441ux43bux430-ux437ux430ux440ux430ux436ux435ux43dux43dux44bux445}
\begin{figure}[H]

{\centering \includegraphics[width=0.7\linewidth,height=\textheight,keepaspectratio]{image/6.png}

}

\caption{Динамика числа зараженных ( I(t) )}

\end{figure}%
\end{block}

\begin{block}{Экспоненциальный рост}
\phantomsection\label{ux44dux43aux441ux43fux43eux43dux435ux43dux446ux438ux430ux43bux44cux43dux44bux439-ux440ux43eux441ux442}
\begin{figure}[H]

{\centering \includegraphics[width=0.7\linewidth,height=\textheight,keepaspectratio]{image/7.png}

}

\caption{Экспоненциальный рост числа зараженных в логарифмическом
масштабе}

\end{figure}%
\end{block}

\begin{block}{Динамика всех групп}
\phantomsection\label{ux434ux438ux43dux430ux43cux438ux43aux430-ux432ux441ux435ux445-ux433ux440ux443ux43fux43f}
\begin{figure}[H]

{\centering \includegraphics[width=0.7\linewidth,height=\textheight,keepaspectratio]{image/8.png}

}

\caption{Модель SIR: динамика всех трёх групп населения}

\end{figure}%
\end{block}

\begin{block}{Компактная панель графиков}
\phantomsection\label{ux43aux43eux43cux43fux430ux43aux442ux43dux430ux44f-ux43fux430ux43dux435ux43bux44c-ux433ux440ux430ux444ux438ux43aux43eux432}
\begin{figure}[H]

{\centering \includegraphics[width=0.7\linewidth,height=\textheight,keepaspectratio]{image/9.png}

}

\caption{Компактная панель графиков для анализа модели SIR}

\end{figure}%
\end{block}

\begin{block}{Динамика в процентах}
\phantomsection\label{ux434ux438ux43dux430ux43cux438ux43aux430-ux432-ux43fux440ux43eux446ux435ux43dux442ux430ux445}
\begin{figure}[H]

{\centering \includegraphics[width=0.7\linewidth,height=\textheight,keepaspectratio]{image/10.png}

}

\caption{Динамика эпидемии в процентах от общей численности популяции}

\end{figure}%
\end{block}

\begin{block}{Фазовый портрет}
\phantomsection\label{ux444ux430ux437ux43eux432ux44bux439-ux43fux43eux440ux442ux440ux435ux442}
\begin{figure}[H]

{\centering \includegraphics[width=0.7\linewidth,height=\textheight,keepaspectratio]{image/11.png}

}

\caption{Фазовый портрет модели SIR}

\end{figure}%
\end{block}
\end{frame}

\begin{frame}{Литературное программирование SIR}
\phantomsection\label{ux43bux438ux442ux435ux440ux430ux442ux443ux440ux43dux43eux435-ux43fux440ux43eux433ux440ux430ux43cux43cux438ux440ux43eux432ux430ux43dux438ux435-sir}
\begin{block}{Литературный код}
\phantomsection\label{ux43bux438ux442ux435ux440ux430ux442ux443ux440ux43dux44bux439-ux43aux43eux434}
\begin{figure}[H]

{\centering \includegraphics[width=0.7\linewidth,height=\textheight,keepaspectratio]{image/12.png}

}

\caption{Литературный код модели SIR в файле sir\_literode.jl}

\end{figure}%
\end{block}

\begin{block}{Генерация форматов}
\phantomsection\label{ux433ux435ux43dux435ux440ux430ux446ux438ux44f-ux444ux43eux440ux43cux430ux442ux43eux432}
\begin{figure}[H]

{\centering \includegraphics[width=0.7\linewidth,height=\textheight,keepaspectratio]{image/13.png}

}

\caption{Запуск скрипта tangle.jl для генерации форматов}

\end{figure}%
\end{block}

\begin{block}{Jupyter Notebook}
\phantomsection\label{jupyter-notebook}
\begin{figure}[H]

{\centering \includegraphics[width=0.7\linewidth,height=\textheight,keepaspectratio]{image/14.png}

}

\caption{Запуск сгенерированного Jupyter Notebook}

\end{figure}%%
\begin{figure}[H]

{\centering \includegraphics[width=0.7\linewidth,height=\textheight,keepaspectratio]{image/15.png}

}

\caption{Содержимое sir\_literode.ipynb в Jupyter Notebook}

\end{figure}%
\end{block}

\begin{block}{Quarto-документация}
\phantomsection\label{quarto-ux434ux43eux43aux443ux43cux435ux43dux442ux430ux446ux438ux44f}
\begin{figure}[H]

{\centering 
\includegraphics[width=0.7\linewidth,height=\textheight,keepaspectratio]{image/16.png}

}

\caption{Содержание файла sir\_literode.qmd}

\end{figure}%
\end{block}

\begin{block}{Чистый код}
\phantomsection\label{ux447ux438ux441ux442ux44bux439-ux43aux43eux434}
\begin{figure}[H]

{\centering \includegraphics[width=0.7\linewidth,height=\textheight,keepaspectratio]{image/17.png}

}

\caption{Содержание файла с чистым кодом sir\_literode.jl}

\end{figure}%
\end{block}
\end{frame}

\begin{frame}{Модель Лотки-Вольтерры}
\phantomsection\label{ux43cux43eux434ux435ux43bux44c-ux43bux43eux442ux43aux438-ux432ux43eux43bux44cux442ux435ux440ux440ux44b-1}
\begin{block}{Создание файла}
\phantomsection\label{ux441ux43eux437ux434ux430ux43dux438ux435-ux444ux430ux439ux43bux430-1}
\begin{figure}[H]

{\centering 
\includegraphics[width=0.7\linewidth,height=\textheight,keepaspectratio]{image/18.png}

}

\caption{Создание файла lv\_ode.jl}

\end{figure}%
\end{block}

\begin{block}{Программный код}
\phantomsection\label{ux43fux440ux43eux433ux440ux430ux43cux43cux43dux44bux439-ux43aux43eux434}
\begin{figure}[H]

{\centering \includegraphics[width=0.7\linewidth,height=\textheight,keepaspectratio]{image/19.png}

}

\caption{Программный код модели Лотки-Вольтерры}

\end{figure}%
\end{block}

\begin{block}{Запуск скрипта}
\phantomsection\label{ux437ux430ux43fux443ux441ux43a-ux441ux43aux440ux438ux43fux442ux430-1}
\begin{figure}[H]

{\centering \includegraphics[width=0.7\linewidth,height=\textheight,keepaspectratio]{image/20.png}

}

\caption{Запуск скрипта lv\_ode.jl}

\end{figure}%
\end{block}

\begin{block}{Результаты в консоли}
\phantomsection\label{ux440ux435ux437ux443ux43bux44cux442ux430ux442ux44b-ux432-ux43aux43eux43dux441ux43eux43bux438}
\begin{figure}[H]

{\centering \includegraphics[width=0.7\linewidth,height=\textheight,keepaspectratio]{image/21.png}

}

\caption{Консольный вывод результатов моделирования}

\end{figure}%
\end{block}
\end{frame}

\begin{frame}{Результаты модели Лотки-Вольтерры}
\phantomsection\label{ux440ux435ux437ux443ux43bux44cux442ux430ux442ux44b-ux43cux43eux434ux435ux43bux438-ux43bux43eux442ux43aux438-ux432ux43eux43bux44cux442ux435ux440ux440ux44b}
\begin{block}{Производные популяций}
\phantomsection\label{ux43fux440ux43eux438ux437ux432ux43eux434ux43dux44bux435-ux43fux43eux43fux443ux43bux44fux446ux438ux439}
\begin{figure}[H]

{\centering \includegraphics[width=0.7\linewidth,height=\textheight,keepaspectratio]{image/22.png}

}

\caption{Производные популяций ( dx/dt ) и ( dy/dt )}

\end{figure}%
\end{block}

\begin{block}{Динамика популяций}
\phantomsection\label{ux434ux438ux43dux430ux43cux438ux43aux430-ux43fux43eux43fux443ux43bux44fux446ux438ux439}
\begin{figure}[H]

{\centering 
\includegraphics[width=0.7\linewidth,height=\textheight,keepaspectratio]{image/23.png}

}

\caption{Динамика популяций в модели Лотки-Вольтерры}

\end{figure}%
\end{block}

\begin{block}{Таблица результатов}
\phantomsection\label{ux442ux430ux431ux43bux438ux446ux430-ux440ux435ux437ux443ux43bux44cux442ux430ux442ux43eux432}
\begin{figure}[H]

{\centering 
\includegraphics[width=0.7\linewidth,height=\textheight,keepaspectratio]{image/24.png}

}

\caption{Таблица с результатами моделирования}

\end{figure}%
\end{block}

\begin{block}{Фазовый портрет}
\phantomsection\label{ux444ux430ux437ux43eux432ux44bux439-ux43fux43eux440ux442ux440ux435ux442-1}
\begin{figure}[H]

{\centering \includegraphics[width=0.7\linewidth,height=\textheight,keepaspectratio]{image/25.png}

}

\caption{Фазовый портрет системы «хищник-жертва»}

\end{figure}%
\end{block}

\begin{block}{Относительные темпы роста}
\phantomsection\label{ux43eux442ux43dux43eux441ux438ux442ux435ux43bux44cux43dux44bux435-ux442ux435ux43cux43fux44b-ux440ux43eux441ux442ux430}
\begin{figure}[H]

{\centering \includegraphics[width=0.7\linewidth,height=\textheight,keepaspectratio]{image/26.png}

}

\caption{Относительные темпы роста популяций}

\end{figure}%
\end{block}

\begin{block}{Спектральный анализ}
\phantomsection\label{ux441ux43fux435ux43aux442ux440ux430ux43bux44cux43dux44bux439-ux430ux43dux430ux43bux438ux437}
\begin{figure}[H]

{\centering \includegraphics[width=0.7\linewidth,height=\textheight,keepaspectratio]{image/27.png}

}

\caption{Спектральный анализ (Фурье) колебаний}

\end{figure}%
\end{block}
\end{frame}

\begin{frame}{Литературное программирование Лотки-Вольтерры}
\phantomsection\label{ux43bux438ux442ux435ux440ux430ux442ux443ux440ux43dux43eux435-ux43fux440ux43eux433ux440ux430ux43cux43cux438ux440ux43eux432ux430ux43dux438ux435-ux43bux43eux442ux43aux438-ux432ux43eux43bux44cux442ux435ux440ux440ux44b}
\begin{block}{Литературный код (часть 1)}
\phantomsection\label{ux43bux438ux442ux435ux440ux430ux442ux443ux440ux43dux44bux439-ux43aux43eux434-ux447ux430ux441ux442ux44c-1}
\begin{figure}[H]

{\centering \includegraphics[width=0.7\linewidth,height=\textheight,keepaspectratio]{image/28.png}

}

\caption{Фрагмент литературного кода (версия 1)}

\end{figure}%
\end{block}

\begin{block}{Литературный код (часть 2)}
\phantomsection\label{ux43bux438ux442ux435ux440ux430ux442ux443ux440ux43dux44bux439-ux43aux43eux434-ux447ux430ux441ux442ux44c-2}
\begin{figure}[H]

{\centering \includegraphics[width=0.7\linewidth,height=\textheight,keepaspectratio]{image/29.png}

}

\caption{Фрагмент литературного кода (версия 2)}

\end{figure}%
\end{block}

\begin{block}{Литературный код (часть 3)}
\phantomsection\label{ux43bux438ux442ux435ux440ux430ux442ux443ux440ux43dux44bux439-ux43aux43eux434-ux447ux430ux441ux442ux44c-3}
\begin{figure}[H]

{\centering \includegraphics[width=0.7\linewidth,height=\textheight,keepaspectratio]{image/30.png}

}

\caption{Фрагмент литературного кода (версия 3)}

\end{figure}%
\end{block}

\begin{block}{Генерация форматов}
\phantomsection\label{ux433ux435ux43dux435ux440ux430ux446ux438ux44f-ux444ux43eux440ux43cux430ux442ux43eux432-1}
\begin{figure}[H]

{\centering \includegraphics[width=0.7\linewidth,height=\textheight,keepaspectratio]{image/31.png}

}

\caption{Генерация производных форматов}

\end{figure}%
\end{block}

\begin{block}{Jupyter Notebook}
\phantomsection\label{jupyter-notebook-1}
\begin{figure}[H]

{\centering \includegraphics[width=0.7\linewidth,height=\textheight,keepaspectratio]{image/32.png}

}

\caption{Запуск Jupyter Notebook для модели Лотки-Вольтерры}

\end{figure}%%
\begin{figure}[H]

{\centering \includegraphics[width=0.7\linewidth,height=\textheight,keepaspectratio]{image/33.png}

}

\caption{Содержимое файла lv\_literate.ipynb}

\end{figure}%
\end{block}
\end{frame}

\section{Результаты}\label{ux440ux435ux437ux443ux43bux44cux442ux430ux442ux44b}

\begin{frame}{Результаты модели SIR}
\phantomsection\label{ux440ux435ux437ux443ux43bux44cux442ux430ux442ux44b-ux43cux43eux434ux435ux43bux438-sir-1}
\begin{itemize}[<+->]
\tightlist
\item
  Подтверждён экспоненциальный рост на начальном этапе
\item
  Пик эпидемии наступает при ( R\_e = 1 )
\item
  Порог коллективного иммунитета: ( 1-1/R\_0 )
\item
  Получены графики динамики всех групп населения
\item
  Построен фазовый портрет системы
\end{itemize}
\end{frame}

\begin{frame}{Результаты модели Лотки-Вольтерры}
\phantomsection\label{ux440ux435ux437ux443ux43bux44cux442ux430ux442ux44b-ux43cux43eux434ux435ux43bux438-ux43bux43eux442ux43aux438-ux432ux43eux43bux44cux442ux435ux440ux440ux44b-1}
\begin{itemize}[<+->]
\tightlist
\item
  Получены циклические решения
\item
  Стационарные точки совпадают с теорией: ( x\^{}* = \gamma/\delta,
  y\^{}* = \alpha/\beta )
\item
  Период колебаний: около 4 единиц времени
\item
  Подтверждён сдвиг фаз между популяциями
\item
  Проведён спектральный анализ колебаний
\end{itemize}
\end{frame}

\begin{frame}{Литературное программирование}
\phantomsection\label{ux43bux438ux442ux435ux440ux430ux442ux443ux440ux43dux43eux435-ux43fux440ux43eux433ux440ux430ux43cux43cux438ux440ux43eux432ux430ux43dux438ux435}
\begin{itemize}[<+->]
\tightlist
\item
  Код для обеих моделей оформлен в литературном стиле с использованием
  Literate.jl
\item
  Сгенерированы:

  \begin{itemize}[<+->]
  \tightlist
  \item
    Чистый код (.jl)
  \item
    Jupyter Notebook (.ipynb)
  \item
    Quarto-документация (.qmd)
  \end{itemize}
\item
  Обеспечена воспроизводимость исследований
\end{itemize}
\end{frame}

\section{Выводы}\label{ux432ux44bux432ux43eux434ux44b}

\begin{frame}{Выводы}
\phantomsection\label{ux432ux44bux432ux43eux434ux44b-1}
\begin{itemize}[<+->]
\tightlist
\item
  Реализованы две классические модели: SIR и Лотки-Вольтерры
\item
  Проведён анализ динамики систем с визуализацией результатов
\item
  Код оформлен в литературном стиле с использованием Literate.jl
\item
  Сгенерированы Jupyter Notebook и Quarto-документация
\item
  Все поставленные задачи выполнены
\item
  Полученные результаты полностью соответствуют теоретическим
  предсказаниям
\end{itemize}
\end{frame}

\begin{frame}{Итоговый слайд}
\phantomsection\label{ux438ux442ux43eux433ux43eux432ux44bux439-ux441ux43bux430ux439ux434}
\begin{itemize}[<+->]
\tightlist
\item
  Моделирование --- мощный инструмент исследования биологических и
  эпидемиологических процессов
\item
  Julia и DrWatson обеспечивают воспроизводимость научных исследований
\item
  Литературное программирование улучшает документирование и понимание
  кода
\item
  Полученные навыки могут быть применены для моделирования других
  процессов
\end{itemize}
\end{frame}




\end{document}
